% =====================================================================
% Cross-corpus zero-shot generalization (DAIC-WOZ -> E-DAIC). Drop into
% paper.tex with % =====================================================================
% Cross-corpus zero-shot generalization (DAIC-WOZ -> E-DAIC). Drop into
% paper.tex with % =====================================================================
% Cross-corpus zero-shot generalization (DAIC-WOZ -> E-DAIC). Drop into
% paper.tex with % =====================================================================
% Cross-corpus zero-shot generalization (DAIC-WOZ -> E-DAIC). Drop into
% paper.tex with \input{sec_edaic}, placed AFTER the fine-tuning section
% (\label{sec:finetune}) and BEFORE \section{Discussion and Limitations}.
% Requires \cite{ringeval2019avec} in the bibliography.
% =====================================================================
\section{Cross-Corpus Generalization: Zero-Shot Transfer to E-DAIC}
\label{sec:edaic}

Every result so far is within DAIC-WOZ: train, select, and test on disjoint
splits of one corpus. The sterner question for a clinical screen is whether the
model transfers to a \emph{different} corpus it never saw --- a different cohort,
a different ASR front end, a different recording protocol. We test this directly
and \emph{zero-shot}: the DAIC-WOZ model is frozen and run, with no adaptation
and no target-side threshold tuning, on E-DAIC \cite{ringeval2019avec}, the
AVEC~2019 successor corpus.

\subsection{Setup}
\label{sec:edaic-setup}
E-DAIC ships automatic-speech-recognition transcripts with \emph{no speaker
diarization}: the virtual interviewer's prompts are not transcribed, so every
segment is participant speech. To make transfer a measurement of \emph{domain}
shift rather than of \emph{input-format} shift, we match instance construction
across corpora: the DAIC-WOZ source model is trained on participant-only chunks
built by silence-gap segmentation (gap $=2.0$\,s; \Cref{sec:method}), exactly the
construction forced by E-DAIC's un-diarized transcripts. We pull the 275 E-DAIC
subjects (train 163 / dev 56 / test 56) and evaluate on the official test split.
The source is the 30-seed bge-large ensemble; its averaged probability is scored
on E-DAIC with \emph{no} refitting. Because the two corpora overlap in
participant-ID ranges, E-DAIC embeddings are cached in a separate namespace to
preclude any cross-corpus reuse.

\textbf{Transfer model fixed a priori.} The configuration carried to E-DAIC was
\emph{not} chosen on the target. The transfer model is the order-free
(temporal-\texttt{none}) pooler, inherited unchanged from the DAIC-WOZ
single-model selection (\Cref{sec:internal}), which had already settled on
order-free pooling for the headline before any E-DAIC data was touched; it was
evaluated on E-DAIC exactly once. The BiGRU variant compared below was
\emph{not} part of that a-priori choice: it surfaced only afterwards, as the
winner of a separate participant-only \emph{development}-AUC ablation on
DAIC-WOZ, and is reported here as a post-hoc probe of whether the in-domain
winner also transfers. Neither configuration's E-DAIC test score was used to
select it --- the order-free model was pre-specified, the BiGRU was dev-selected
on the source corpus --- so the comparison below describes, and does not
circularly choose, the better generalizer.

We report the threshold-free ROC-AUC and PR-AUC as the primary transfer numbers:
a decision threshold cannot be tuned on the target without leaking it, and the
source threshold is mis-calibrated under domain shift (below).

\subsection{Result: the model transfers, and order-free pooling transfers best}
\label{sec:edaic-result}

\Cref{tab:edaic} reports in-domain DAIC-WOZ alongside zero-shot E-DAIC for the two
temporal-pooling variants of the matched source. Two findings.

\textbf{(1) The model generalizes, with the expected domain-shift cost.} The
order-free model carries an in-domain test AUC of $0.819$ to a zero-shot E-DAIC
AUC of $0.716$ (30-seed per-seed mean $0.716\pm0.017$; subject-level bootstrap
$95\%$ CI $[.55,.88]$; PR-AUC $0.685$) --- a $\approx\!0.10$ AUC drop, well above
chance and consistent with the cross-ASR, cross-cohort gap. As anticipated, the thresholded operating point degrades more
than the ranking: with the source threshold transferred unchanged the F1 is
$0.46$, and even a target-prevalence threshold reaches only F1~$0.54$ /
balanced-accuracy~$0.67$. Domain shift de-calibrates the probabilities while
leaving the ranking largely intact, which is why AUC is the honest transfer
metric.

\textbf{(2) The best in-domain configuration is not the best transfer
configuration.} Adding a BiGRU temporal-context layer is the strongest source
model \emph{on DAIC-WOZ development} (per-seed dev AUC $0.900$ vs.\ $0.857$ for
order-free pooling). Yet it transfers \emph{worse}: zero-shot E-DAIC AUC
$0.662$ vs.\ $0.716$, a $0.052$ gap that is significant under a seed-paired
Wilcoxon signed-rank test over the 30 seeds ($p<10^{-4}$). The GRU also fails to
convert its development edge into a DAIC \emph{test} gain ($0.786$ vs.\ $0.819$),
so its advantage is confined to the development split it was selected on: the
temporal layer fits corpus-specific turn-ordering that does not survive a change
of corpus, whereas order-free attention pooling is more robust to the shift.
Selecting a transfer model on in-domain development AUC would therefore pick the
\emph{wrong} model; the order-free pooler is both the safer in-domain choice and
the better generalizer.

\begin{table}[t]
\caption{Zero-shot cross-corpus transfer. Source: bge-large, participant-only,
gap-merged, trained on DAIC-WOZ; evaluated on E-DAIC with no target adaptation.
ROC-AUC as 30-seed per-seed mean$\pm$std; the E-DAIC column adds a subject-level
bootstrap 95\% CI (2000 resamples of the 56 test subjects) on the seed-ensemble.
The order-free model wins the transfer despite the GRU winning the DAIC
development split.}
\label{tab:edaic}
\centering\small
\setlength{\tabcolsep}{4pt}
\resizebox{\columnwidth}{!}{%
\begin{tabular}{@{}lccc@{}}
\toprule
& \multicolumn{2}{c}{\textbf{DAIC-WOZ (in-domain)}} & \textbf{E-DAIC (zero-shot)} \\
\cmidrule(lr){2-3}\cmidrule(lr){4-4}
\textbf{Source head} & \textbf{dev AUC} & \textbf{test AUC} & \textbf{test AUC} \\
\midrule
order-free (none) & 0.857\,{\footnotesize$\pm$.014} & \textbf{0.819}\,{\footnotesize$\pm$.019} & \textbf{0.716}\,{\footnotesize$\pm$.017\,[.55,.88]} \\
BiGRU             & \textbf{0.900}\,{\footnotesize$\pm$.007} & 0.786\,{\footnotesize$\pm$.040} & 0.662\,{\footnotesize$\pm$.011\,[.48,.83]} \\
\bottomrule
\end{tabular}}
\end{table}

\subsection{Takeaway}
\label{sec:edaic-takeaway}
A text-only DAIC-WOZ model transfers zero-shot to E-DAIC at AUC $\approx\!0.72$
with no target supervision --- evidence that the gated-attention MIL screen
learns corpus-transferable depression signal rather than a single corpus's
artefacts. The transfer also exposes a selection lesson that the rest of this
paper's in-domain protocol cannot: a temporal-context head that helps in-domain
\emph{hurts} out-of-corpus, so generalization-critical design choices should be
validated across corpora, not on a single development split.
, placed AFTER the fine-tuning section
% (\label{sec:finetune}) and BEFORE \section{Discussion and Limitations}.
% Requires \cite{ringeval2019avec} in the bibliography.
% =====================================================================
\section{Cross-Corpus Generalization: Zero-Shot Transfer to E-DAIC}
\label{sec:edaic}

Every result so far is within DAIC-WOZ: train, select, and test on disjoint
splits of one corpus. The sterner question for a clinical screen is whether the
model transfers to a \emph{different} corpus it never saw --- a different cohort,
a different ASR front end, a different recording protocol. We test this directly
and \emph{zero-shot}: the DAIC-WOZ model is frozen and run, with no adaptation
and no target-side threshold tuning, on E-DAIC \cite{ringeval2019avec}, the
AVEC~2019 successor corpus.

\subsection{Setup}
\label{sec:edaic-setup}
E-DAIC ships automatic-speech-recognition transcripts with \emph{no speaker
diarization}: the virtual interviewer's prompts are not transcribed, so every
segment is participant speech. To make transfer a measurement of \emph{domain}
shift rather than of \emph{input-format} shift, we match instance construction
across corpora: the DAIC-WOZ source model is trained on participant-only chunks
built by silence-gap segmentation (gap $=2.0$\,s; \Cref{sec:method}), exactly the
construction forced by E-DAIC's un-diarized transcripts. We pull the 275 E-DAIC
subjects (train 163 / dev 56 / test 56) and evaluate on the official test split.
The source is the 30-seed bge-large ensemble; its averaged probability is scored
on E-DAIC with \emph{no} refitting. Because the two corpora overlap in
participant-ID ranges, E-DAIC embeddings are cached in a separate namespace to
preclude any cross-corpus reuse.

\textbf{Transfer model fixed a priori.} The configuration carried to E-DAIC was
\emph{not} chosen on the target. The transfer model is the order-free
(temporal-\texttt{none}) pooler, inherited unchanged from the DAIC-WOZ
single-model selection (\Cref{sec:internal}), which had already settled on
order-free pooling for the headline before any E-DAIC data was touched; it was
evaluated on E-DAIC exactly once. The BiGRU variant compared below was
\emph{not} part of that a-priori choice: it surfaced only afterwards, as the
winner of a separate participant-only \emph{development}-AUC ablation on
DAIC-WOZ, and is reported here as a post-hoc probe of whether the in-domain
winner also transfers. Neither configuration's E-DAIC test score was used to
select it --- the order-free model was pre-specified, the BiGRU was dev-selected
on the source corpus --- so the comparison below describes, and does not
circularly choose, the better generalizer.

We report the threshold-free ROC-AUC and PR-AUC as the primary transfer numbers:
a decision threshold cannot be tuned on the target without leaking it, and the
source threshold is mis-calibrated under domain shift (below).

\subsection{Result: the model transfers, and order-free pooling transfers best}
\label{sec:edaic-result}

\Cref{tab:edaic} reports in-domain DAIC-WOZ alongside zero-shot E-DAIC for the two
temporal-pooling variants of the matched source. Two findings.

\textbf{(1) The model generalizes, with the expected domain-shift cost.} The
order-free model carries an in-domain test AUC of $0.819$ to a zero-shot E-DAIC
AUC of $0.716$ (30-seed per-seed mean $0.716\pm0.017$; subject-level bootstrap
$95\%$ CI $[.55,.88]$; PR-AUC $0.685$) --- a $\approx\!0.10$ AUC drop, well above
chance and consistent with the cross-ASR, cross-cohort gap. As anticipated, the thresholded operating point degrades more
than the ranking: with the source threshold transferred unchanged the F1 is
$0.46$, and even a target-prevalence threshold reaches only F1~$0.54$ /
balanced-accuracy~$0.67$. Domain shift de-calibrates the probabilities while
leaving the ranking largely intact, which is why AUC is the honest transfer
metric.

\textbf{(2) The best in-domain configuration is not the best transfer
configuration.} Adding a BiGRU temporal-context layer is the strongest source
model \emph{on DAIC-WOZ development} (per-seed dev AUC $0.900$ vs.\ $0.857$ for
order-free pooling). Yet it transfers \emph{worse}: zero-shot E-DAIC AUC
$0.662$ vs.\ $0.716$, a $0.052$ gap that is significant under a seed-paired
Wilcoxon signed-rank test over the 30 seeds ($p<10^{-4}$). The GRU also fails to
convert its development edge into a DAIC \emph{test} gain ($0.786$ vs.\ $0.819$),
so its advantage is confined to the development split it was selected on: the
temporal layer fits corpus-specific turn-ordering that does not survive a change
of corpus, whereas order-free attention pooling is more robust to the shift.
Selecting a transfer model on in-domain development AUC would therefore pick the
\emph{wrong} model; the order-free pooler is both the safer in-domain choice and
the better generalizer.

\begin{table}[t]
\caption{Zero-shot cross-corpus transfer. Source: bge-large, participant-only,
gap-merged, trained on DAIC-WOZ; evaluated on E-DAIC with no target adaptation.
ROC-AUC as 30-seed per-seed mean$\pm$std; the E-DAIC column adds a subject-level
bootstrap 95\% CI (2000 resamples of the 56 test subjects) on the seed-ensemble.
The order-free model wins the transfer despite the GRU winning the DAIC
development split.}
\label{tab:edaic}
\centering\small
\setlength{\tabcolsep}{4pt}
\resizebox{\columnwidth}{!}{%
\begin{tabular}{@{}lccc@{}}
\toprule
& \multicolumn{2}{c}{\textbf{DAIC-WOZ (in-domain)}} & \textbf{E-DAIC (zero-shot)} \\
\cmidrule(lr){2-3}\cmidrule(lr){4-4}
\textbf{Source head} & \textbf{dev AUC} & \textbf{test AUC} & \textbf{test AUC} \\
\midrule
order-free (none) & 0.857\,{\footnotesize$\pm$.014} & \textbf{0.819}\,{\footnotesize$\pm$.019} & \textbf{0.716}\,{\footnotesize$\pm$.017\,[.55,.88]} \\
BiGRU             & \textbf{0.900}\,{\footnotesize$\pm$.007} & 0.786\,{\footnotesize$\pm$.040} & 0.662\,{\footnotesize$\pm$.011\,[.48,.83]} \\
\bottomrule
\end{tabular}}
\end{table}

\subsection{Takeaway}
\label{sec:edaic-takeaway}
A text-only DAIC-WOZ model transfers zero-shot to E-DAIC at AUC $\approx\!0.72$
with no target supervision --- evidence that the gated-attention MIL screen
learns corpus-transferable depression signal rather than a single corpus's
artefacts. The transfer also exposes a selection lesson that the rest of this
paper's in-domain protocol cannot: a temporal-context head that helps in-domain
\emph{hurts} out-of-corpus, so generalization-critical design choices should be
validated across corpora, not on a single development split.
, placed AFTER the fine-tuning section
% (\label{sec:finetune}) and BEFORE \section{Discussion and Limitations}.
% Requires \cite{ringeval2019avec} in the bibliography.
% =====================================================================
\section{Cross-Corpus Generalization: Zero-Shot Transfer to E-DAIC}
\label{sec:edaic}

Every result so far is within DAIC-WOZ: train, select, and test on disjoint
splits of one corpus. The sterner question for a clinical screen is whether the
model transfers to a \emph{different} corpus it never saw --- a different cohort,
a different ASR front end, a different recording protocol. We test this directly
and \emph{zero-shot}: the DAIC-WOZ model is frozen and run, with no adaptation
and no target-side threshold tuning, on E-DAIC \cite{ringeval2019avec}, the
AVEC~2019 successor corpus.

\subsection{Setup}
\label{sec:edaic-setup}
E-DAIC ships automatic-speech-recognition transcripts with \emph{no speaker
diarization}: the virtual interviewer's prompts are not transcribed, so every
segment is participant speech. To make transfer a measurement of \emph{domain}
shift rather than of \emph{input-format} shift, we match instance construction
across corpora: the DAIC-WOZ source model is trained on participant-only chunks
built by silence-gap segmentation (gap $=2.0$\,s; \Cref{sec:method}), exactly the
construction forced by E-DAIC's un-diarized transcripts. We pull the 275 E-DAIC
subjects (train 163 / dev 56 / test 56) and evaluate on the official test split.
The source is the 30-seed bge-large ensemble; its averaged probability is scored
on E-DAIC with \emph{no} refitting. Because the two corpora overlap in
participant-ID ranges, E-DAIC embeddings are cached in a separate namespace to
preclude any cross-corpus reuse.

\textbf{Transfer model fixed a priori.} The configuration carried to E-DAIC was
\emph{not} chosen on the target. The transfer model is the order-free
(temporal-\texttt{none}) pooler, inherited unchanged from the DAIC-WOZ
single-model selection (\Cref{sec:internal}), which had already settled on
order-free pooling for the headline before any E-DAIC data was touched; it was
evaluated on E-DAIC exactly once. The BiGRU variant compared below was
\emph{not} part of that a-priori choice: it surfaced only afterwards, as the
winner of a separate participant-only \emph{development}-AUC ablation on
DAIC-WOZ, and is reported here as a post-hoc probe of whether the in-domain
winner also transfers. Neither configuration's E-DAIC test score was used to
select it --- the order-free model was pre-specified, the BiGRU was dev-selected
on the source corpus --- so the comparison below describes, and does not
circularly choose, the better generalizer.

We report the threshold-free ROC-AUC and PR-AUC as the primary transfer numbers:
a decision threshold cannot be tuned on the target without leaking it, and the
source threshold is mis-calibrated under domain shift (below).

\subsection{Result: the model transfers, and order-free pooling transfers best}
\label{sec:edaic-result}

\Cref{tab:edaic} reports in-domain DAIC-WOZ alongside zero-shot E-DAIC for the two
temporal-pooling variants of the matched source. Two findings.

\textbf{(1) The model generalizes, with the expected domain-shift cost.} The
order-free model carries an in-domain test AUC of $0.819$ to a zero-shot E-DAIC
AUC of $0.716$ (30-seed per-seed mean $0.716\pm0.017$; subject-level bootstrap
$95\%$ CI $[.55,.88]$; PR-AUC $0.685$) --- a $\approx\!0.10$ AUC drop, well above
chance and consistent with the cross-ASR, cross-cohort gap. As anticipated, the thresholded operating point degrades more
than the ranking: with the source threshold transferred unchanged the F1 is
$0.46$, and even a target-prevalence threshold reaches only F1~$0.54$ /
balanced-accuracy~$0.67$. Domain shift de-calibrates the probabilities while
leaving the ranking largely intact, which is why AUC is the honest transfer
metric.

\textbf{(2) The best in-domain configuration is not the best transfer
configuration.} Adding a BiGRU temporal-context layer is the strongest source
model \emph{on DAIC-WOZ development} (per-seed dev AUC $0.900$ vs.\ $0.857$ for
order-free pooling). Yet it transfers \emph{worse}: zero-shot E-DAIC AUC
$0.662$ vs.\ $0.716$, a $0.052$ gap that is significant under a seed-paired
Wilcoxon signed-rank test over the 30 seeds ($p<10^{-4}$). The GRU also fails to
convert its development edge into a DAIC \emph{test} gain ($0.786$ vs.\ $0.819$),
so its advantage is confined to the development split it was selected on: the
temporal layer fits corpus-specific turn-ordering that does not survive a change
of corpus, whereas order-free attention pooling is more robust to the shift.
Selecting a transfer model on in-domain development AUC would therefore pick the
\emph{wrong} model; the order-free pooler is both the safer in-domain choice and
the better generalizer.

\begin{table}[t]
\caption{Zero-shot cross-corpus transfer. Source: bge-large, participant-only,
gap-merged, trained on DAIC-WOZ; evaluated on E-DAIC with no target adaptation.
ROC-AUC as 30-seed per-seed mean$\pm$std; the E-DAIC column adds a subject-level
bootstrap 95\% CI (2000 resamples of the 56 test subjects) on the seed-ensemble.
The order-free model wins the transfer despite the GRU winning the DAIC
development split.}
\label{tab:edaic}
\centering\small
\setlength{\tabcolsep}{4pt}
\resizebox{\columnwidth}{!}{%
\begin{tabular}{@{}lccc@{}}
\toprule
& \multicolumn{2}{c}{\textbf{DAIC-WOZ (in-domain)}} & \textbf{E-DAIC (zero-shot)} \\
\cmidrule(lr){2-3}\cmidrule(lr){4-4}
\textbf{Source head} & \textbf{dev AUC} & \textbf{test AUC} & \textbf{test AUC} \\
\midrule
order-free (none) & 0.857\,{\footnotesize$\pm$.014} & \textbf{0.819}\,{\footnotesize$\pm$.019} & \textbf{0.716}\,{\footnotesize$\pm$.017\,[.55,.88]} \\
BiGRU             & \textbf{0.900}\,{\footnotesize$\pm$.007} & 0.786\,{\footnotesize$\pm$.040} & 0.662\,{\footnotesize$\pm$.011\,[.48,.83]} \\
\bottomrule
\end{tabular}}
\end{table}

\subsection{Takeaway}
\label{sec:edaic-takeaway}
A text-only DAIC-WOZ model transfers zero-shot to E-DAIC at AUC $\approx\!0.72$
with no target supervision --- evidence that the gated-attention MIL screen
learns corpus-transferable depression signal rather than a single corpus's
artefacts. The transfer also exposes a selection lesson that the rest of this
paper's in-domain protocol cannot: a temporal-context head that helps in-domain
\emph{hurts} out-of-corpus, so generalization-critical design choices should be
validated across corpora, not on a single development split.
, placed AFTER the fine-tuning section
% (\label{sec:finetune}) and BEFORE \section{Discussion and Limitations}.
% Requires \cite{ringeval2019avec} in the bibliography.
% =====================================================================
\section{Cross-Corpus Generalization: Zero-Shot Transfer to E-DAIC}
\label{sec:edaic}

Every result so far is within DAIC-WOZ: train, select, and test on disjoint
splits of one corpus. The sterner question for a clinical screen is whether the
model transfers to a \emph{different} corpus it never saw --- a different cohort,
a different ASR front end, a different recording protocol. We test this directly
and \emph{zero-shot}: the DAIC-WOZ model is frozen and run, with no adaptation
and no target-side threshold tuning, on E-DAIC \cite{ringeval2019avec}, the
AVEC~2019 successor corpus.

\subsection{Setup}
\label{sec:edaic-setup}
E-DAIC ships automatic-speech-recognition transcripts with \emph{no speaker
diarization}: the virtual interviewer's prompts are not transcribed, so every
segment is participant speech. To make transfer a measurement of \emph{domain}
shift rather than of \emph{input-format} shift, we match instance construction
across corpora: the DAIC-WOZ source model is trained on participant-only chunks
built by silence-gap segmentation (gap $=2.0$\,s; \Cref{sec:method}), exactly the
construction forced by E-DAIC's un-diarized transcripts. We pull the 275 E-DAIC
subjects (train 163 / dev 56 / test 56) and evaluate on the official test split.
The source is the 30-seed bge-large ensemble; its averaged probability is scored
on E-DAIC with \emph{no} refitting. Because the two corpora overlap in
participant-ID ranges, E-DAIC embeddings are cached in a separate namespace to
preclude any cross-corpus reuse.

\textbf{Transfer model fixed a priori.} The configuration carried to E-DAIC was
\emph{not} chosen on the target. The transfer model is the order-free
(temporal-\texttt{none}) pooler, inherited unchanged from the DAIC-WOZ
single-model selection (\Cref{sec:internal}), which had already settled on
order-free pooling for the headline before any E-DAIC data was touched; it was
evaluated on E-DAIC exactly once. The BiGRU variant compared below was
\emph{not} part of that a-priori choice: it surfaced only afterwards, as the
winner of a separate participant-only \emph{development}-AUC ablation on
DAIC-WOZ, and is reported here as a post-hoc probe of whether the in-domain
winner also transfers. Neither configuration's E-DAIC test score was used to
select it --- the order-free model was pre-specified, the BiGRU was dev-selected
on the source corpus --- so the comparison below describes, and does not
circularly choose, the better generalizer.

We report the threshold-free ROC-AUC and PR-AUC as the primary transfer numbers:
a decision threshold cannot be tuned on the target without leaking it, and the
source threshold is mis-calibrated under domain shift (below).

\subsection{Result: the model transfers, and order-free pooling transfers best}
\label{sec:edaic-result}

\Cref{tab:edaic} reports in-domain DAIC-WOZ alongside zero-shot E-DAIC for the two
temporal-pooling variants of the matched source. Two findings.

\textbf{(1) The model generalizes, with the expected domain-shift cost.} The
order-free model carries an in-domain test AUC of $0.819$ to a zero-shot E-DAIC
AUC of $0.716$ (30-seed per-seed mean $0.716\pm0.017$; subject-level bootstrap
$95\%$ CI $[.55,.88]$; PR-AUC $0.685$) --- a $\approx\!0.10$ AUC drop, well above
chance and consistent with the cross-ASR, cross-cohort gap. As anticipated, the thresholded operating point degrades more
than the ranking: with the source threshold transferred unchanged the F1 is
$0.46$, and even a target-prevalence threshold reaches only F1~$0.54$ /
balanced-accuracy~$0.67$. Domain shift de-calibrates the probabilities while
leaving the ranking largely intact, which is why AUC is the honest transfer
metric.

\textbf{(2) The best in-domain configuration is not the best transfer
configuration.} Adding a BiGRU temporal-context layer is the strongest source
model \emph{on DAIC-WOZ development} (per-seed dev AUC $0.900$ vs.\ $0.857$ for
order-free pooling). Yet it transfers \emph{worse}: zero-shot E-DAIC AUC
$0.662$ vs.\ $0.716$, a $0.052$ gap that is significant under a seed-paired
Wilcoxon signed-rank test over the 30 seeds ($p<10^{-4}$). The GRU also fails to
convert its development edge into a DAIC \emph{test} gain ($0.786$ vs.\ $0.819$),
so its advantage is confined to the development split it was selected on: the
temporal layer fits corpus-specific turn-ordering that does not survive a change
of corpus, whereas order-free attention pooling is more robust to the shift.
Selecting a transfer model on in-domain development AUC would therefore pick the
\emph{wrong} model; the order-free pooler is both the safer in-domain choice and
the better generalizer.

\begin{table}[t]
\caption{Zero-shot cross-corpus transfer. Source: bge-large, participant-only,
gap-merged, trained on DAIC-WOZ; evaluated on E-DAIC with no target adaptation.
ROC-AUC as 30-seed per-seed mean$\pm$std; the E-DAIC column adds a subject-level
bootstrap 95\% CI (2000 resamples of the 56 test subjects) on the seed-ensemble.
The order-free model wins the transfer despite the GRU winning the DAIC
development split.}
\label{tab:edaic}
\centering\small
\setlength{\tabcolsep}{4pt}
\resizebox{\columnwidth}{!}{%
\begin{tabular}{@{}lccc@{}}
\toprule
& \multicolumn{2}{c}{\textbf{DAIC-WOZ (in-domain)}} & \textbf{E-DAIC (zero-shot)} \\
\cmidrule(lr){2-3}\cmidrule(lr){4-4}
\textbf{Source head} & \textbf{dev AUC} & \textbf{test AUC} & \textbf{test AUC} \\
\midrule
order-free (none) & 0.857\,{\footnotesize$\pm$.014} & \textbf{0.819}\,{\footnotesize$\pm$.019} & \textbf{0.716}\,{\footnotesize$\pm$.017\,[.55,.88]} \\
BiGRU             & \textbf{0.900}\,{\footnotesize$\pm$.007} & 0.786\,{\footnotesize$\pm$.040} & 0.662\,{\footnotesize$\pm$.011\,[.48,.83]} \\
\bottomrule
\end{tabular}}
\end{table}

\subsection{Takeaway}
\label{sec:edaic-takeaway}
A text-only DAIC-WOZ model transfers zero-shot to E-DAIC at AUC $\approx\!0.72$
with no target supervision --- evidence that the gated-attention MIL screen
learns corpus-transferable depression signal rather than a single corpus's
artefacts. The transfer also exposes a selection lesson that the rest of this
paper's in-domain protocol cannot: a temporal-context head that helps in-domain
\emph{hurts} out-of-corpus, so generalization-critical design choices should be
validated across corpora, not on a single development split.
